\documentclass[11pt,a4paper]{article}
\usepackage[T1]{fontenc}
\usepackage{lmodern,amsmath,amssymb,amsthm,microtype}
\usepackage[margin=27mm]{geometry}
\usepackage[colorlinks=true,linkcolor=blue,citecolor=blue,urlcolor=blue]{hyperref}
\setlength{\emergencystretch}{2em}
\newtheorem{proposition}{Proposition}[section]
\newcommand{\Q}{\mathbb Q}
\DeclareMathOperator{\Tr}{Tr}
\title{Odd halfpoint cosets do not automatically prepare raw abscissas}
\author{}\date{17 November 2026}
\begin{document}\maketitle
\begin{abstract}
A proposed raw-abscissa variant of the derivative detector still has
simultaneous poles in its two slots after restricting to an odd input
and odd halfpoint multiples. An exact physical example at the prime
$85817$ has pole-input valuation $-11$. Its actual-trace square
radicand has valuation one. All pre-pole whole-character guards and
real preparation hold for sufficiently large slopes in a specified
arithmetic progression. This disproves automatic preparation for that
variant; it does not disprove completeness using other witnesses.
\end{abstract}

\section{The precise variant and tuple}

Retain the curve, generator, six tests, point algebras and physical
nine-contact map of the derivative construction
\cite{Derivative,Full}, and write
\[
 K=\Q(\sqrt2),\quad E:y^2=x^3-(11-6\sqrt2)x,
 \quad R_0=(1+\sqrt2,2\sqrt2),\quad P=2R_0.
\]
The retained group premise makes $R_0$ non-torsion. In particular,
its nonzero multiples have finite nonzero abscissas and ordinates.
Use the full auxiliary-point domain of \cite{Full}. Consider the
specific replacement
\[
 W_{\rm raw}=\prod_{i=1}^4\left(1+\frac{C_i^2}{u}\right)
       \prod_{i=5}^6\left(1+\frac{D_i^2}{u}\right),
 \qquad D_i=\frac{C_i^2}{1+uC_i^2}\quad(i=5,6).
\]
Put $V=\Tr(W_{\rm raw}^{-1})$, $w=u/(1+u)^2$,
$b=4w/V^2$, and $\mathcal C=V^2+w(2-V)$.
The first four tests are still reciprocal abscissas; only their
balancing factors have been removed. This is the raw squared
variant, not the unsquared balancing proved in \cite{Even}.

Take
\[
 X=1,\quad B=3,\quad A=P,\quad H=U=2665R_0.
\]
Thus $X$ is odd and both auxiliaries lie in the odd halfpoint coset.
Since $U=H$ and $U-A=H-P$, the selected rational character has
\[
 (C_1,\ldots,C_6)=(f(H),f(H),\overline{f(H-P)},
                         \overline{f(H-P)},0,0),\qquad f=1/x.
\]
The physical scale and map are
\[
 a=\left(\frac{10621537}{256}\right)^4,\quad
 F=\prod_{c\in\{-10,-8,-5,1,2,3,4,6,7\}}((a+c)^2-z^2),
 \qquad u=-18F/F_a,
\]
where differentiation in $a$ precedes specialization. The choice
$B=3$ is an algebraic fixture, with no assertion concerning the
ineffective defining-shift threshold.

\begin{proposition}\label{prop:counterexample}
Let $p=85817$. For every integer $z\equiv39827\pmod{p^2}$,
the selected rational character satisfies
\[
 v_p(u)=-1,\quad v_p(w)=1,\quad v_p(V)=6,
 \quad v_p(b)=-11,\quad v_p(\mathcal C)=1.
\]
For all sufficiently large positive such $z$, every retained
pre-pole whole-character guard and the real preparation hold.
Nevertheless neither a rational actual-trace square nor rational
native intrinsic-pole witnesses exist for that character.
\end{proposition}

\begin{proof}
The two roots of $2$ modulo $p$ are $21381$ and $64436$.
The reduction orders of $R_0$ in these slots are respectively
$5330=2\cdot2665$ and $1066=2\cdot533$. In both slots,
$[2665]R_0=(0,0)$. The exact computations modulo $p^3$ give
\[
\begin{array}{c|rr}
 \sqrt2\pmod p&21381&64436\\\hline
 x(H)/p^2\pmod p&41859&53375\\
 y(H)/p\pmod p&67362&36131
\end{array}
\]
with each displayed leading residue nonzero. Thus $v(x(H))=2$.
The root point $P$ has odd order greater than one in each slot;
therefore $(0,0)-P$ has nonzero abscissa and ordinate. Consequently
the first four test valuations are $(-2,-2,0,0)$ in both slots.

For the specified slope class, direct differentiation and modular
evaluation give
\[
 F\equiv7357586644\pmod{p^2},\qquad
 F_a\equiv6965422622=81166p\pmod{p^2}.
\]
Here $F$ is a unit and $81166\not\equiv0\pmod p$. Hence
$v_p(u)=-1$ and $pu\equiv65426\pmod p$. Each of the first two
raw factors has valuation $-3$, while the other factors are units.
Therefore $W_{\rm raw}^{-1}$ has valuation six in each slot.
Its trace has leading coefficient
\[
 \frac{V}{p^6}\equiv
 65426^2(41859^4+53375^4)\equiv33198\not\equiv0\pmod p.
\]
The claimed valuations of $V,w,b$ follow. In $\mathcal C$ the
first term has valuation twelve, and the second valuation one
because $2-V$ is a unit. Thus $\mathcal C$ is nonzero and cannot
be a rational square. The native pole implication
\cite[Theorem~1.1]{Intrinsic} excludes its rational witnesses
because $v_p(7b)=-11$.

It remains to justify the whole-character and real assertions.
The rational point-coordinate pairs are the exact coordinates of
the stated multiples; the multiplication specification suffices
without printing their large numerators. At each real embedding,
the point-algebra characters independently change ordinate signs.
All resulting points are signed multiples $\pm2R_0$ and
$\pm2665R_0$. Their required differences have indices
$\pm2665\pm2$, which are nonzero. Non-torsion therefore excludes
zero ordinates, zero abscissas, and equal-abscissa exceptional
differences. The original point, difference and quotient guards
hold at every character.

The finitely many resulting real tests are fixed as $z$ varies.
For $z>a+10$, all factors of $F$ are negative and $F_a/F<0$,
so $u>0$. As $z\longrightarrow+\infty$, $u\sim z^2/a$,
$w\sim a/z^2$, and every raw or balanced factor tends to one
at every character. The comparison denominators
$1+uC_i^2$ are positive, as are the raw factors. Hence all their
norms and both detector slots are nonzero, $V>0$, and eventually
$0<7b<1$ simultaneously at all characters. The remaining scalar
guards follow from $z>a+10$ and the nonzero physical factors.
Every sufficiently large integer in the stated progression meets
these requirements.
\end{proof}

\section{Scope and exact verification}

No selected-prime factorization, complete fixed finite-place
preparation, or completeness theorem for the odd auxiliary coset
is asserted. Indeed $H=hR_0$, $H-P=(h-2)R_0$ with $h$ odd
can lie in the relevant formal groups only where the corresponding
reduction orders of $R_0$ are odd. Any completeness argument for
this coset must verify those selected-contact conditions and must
also control the new negative-weight primes. The proposition shows
that oddness alone does not do the latter. It leaves other tuples
and additional compatible congruences open.

The standard-library checker in
\nolinkurl{negative_results/computation/pole_parameters/odd_coset_raw_v76/}
computes an affine addition chain modulo $p^3$, verifying that
every denominator is a unit, and directly differentiates the contact
product. A producer-free second checker reconstructs the point
coordinates using division polynomials and evaluates $F_a/F$
as a sum of logarithmic derivatives. Both verify the decisive
integer residues, exact reduction orders and trace noncancellation.

\begin{thebibliography}{9}
\bibitem{Derivative}\emph{An ordinary-$abc$ derivative detector in
eleven variables},
\nolinkurl{papers/conditional_few_variables/derivative_detector11_20261117_v63.tex}.
\bibitem{Full}\emph{Full auxiliary elliptic points in the eleven-variable
derivative source},
\nolinkurl{papers/conditional_few_variables/full_auxiliary11_20261117_v71.tex}.
\bibitem{Even}\emph{Even abscissa valuations in the derivative detector},
\nolinkurl{papers/conditional_few_variables/even_valuation_detector_20261117_v75.tex}.
\bibitem{Intrinsic}\emph{An intrinsic two-parameter pole equation},
\nolinkurl{papers/pole_parameters/intrinsic_pole_20261117_v56.tex}.
\end{thebibliography}
\end{document}
